% Section 8: what the data suggest, kept to one page. Lessons from the evidence
% and from what other places have done. Framed as options, not prescriptions.
\section{What the data suggest}
\label{sec:lessons}

The evidence in this study points toward options rather than a program, and
each row below shows the evidence behind its option so a reader can judge the
inference for themselves. None of these is a recommendation about how much
Austin or Texas should spend, which is a political judgment this study does not
make. The table names the source behind each row in the second column and, in
the third, the closest precedent this study located; every finding is developed
in the section listed, and every underlying figure is registered in the numbers
ledger at
\texttt{03\_analysis/out/canonical\_numbers.csv}.

\vspace{2pt}
{\footnotesize
% Tighter intercolumn padding: with three text columns the default 6pt on each
% side of every column consumes about 9mm of the text block.
\setlength{\tabcolsep}{3.5pt}
% tabularx guarantees the table fits the text block exactly, rather than
% relying on fixed widths that have to be re-tuned whenever a cell changes.
\begin{longtable}{@{}>{\bfseries\color{amnavy}\RaggedRight\arraybackslash}p{4.3cm} >{\RaggedRight\arraybackslash}p{6.6cm} >{\RaggedRight\arraybackslash}p{5.7cm}@{}}
\toprule
\textbf{What the evidence shows} & \textbf{Option it points toward} & \textbf{Precedent or supporting evidence} \\
\midrule
\endfirsthead
\toprule
\textbf{What the evidence shows} & \textbf{Option it points toward} & \textbf{Precedent or supporting evidence} \\
\midrule
\endhead
Austin's own music census stopped asking respondents what they earned after its
2014 fielding, whose income questions covered calendar 2013, and no local
survey has replaced the question.\footnote{Titan Music Group, \emph{The Austin
Music Census: A Data-Driven Assessment of Austin's Commercial Music Economy},
published for the City of Austin Economic Development Department, June 2015,
fielded through community mailing lists and industry channels in late 2014,
n=3{,}968 self-selected online respondents, at
\url{https://austin.widen.net/content/om7zqn1qbx/pdf/Austin_Music_Census_Interactive_PDF_53115.pdf};
Sound Music Cities and the City of Austin Music \& Entertainment Division,
\emph{2022 Greater Austin Music Census}, fielded 15 July to 12 September 2022,
n=2{,}260 self-selected online respondents, at
\url{https://www.austinmusiccensus.org/s/2022-Greater-Austin-Music-Census-SUMMARY-REPORT-v13-optimized.pdf}.
The 2022 instrument replaced the 2014 income-quantification items with an
ordinal share-of-income scale that reports no dollars.}
& Consider restoring a short earnings module to the next music survey and
running it on a fixed cycle. Five consistent questions asked every two years
rebuild the series; consistency matters more than length.
& Other cities in the same census cohort still publish musician income
figures; Section~\ref{sec:measurement} lists fifteen of them. The 2022 Austin
instrument dropped the income-quantification items
its 2014 predecessor asked and replaced them with an ordinal share-of-income
scale that reports no dollars. \\
\rowrule
The Austin Live Music Fund received \NCityLmfFyTwoFourApplicants\
applications in its fiscal 2024 cycle and \NCityLmfFyTwoThreeApplicants\ in
fiscal 2023, both figures from press reporting rather than a primary
tabulation, and neither cycle's applicant-level distribution has been
published.\footnote{The applicant counts trace to \emph{Austin Chronicle}
reporting on the fund (FY2024 count published 28 March 2025); the fund's own
FY2023 dashboard offers no static export,
and no primary applicant tabulation was located for either cycle. Development
in Section~\ref{sec:city}. Both cycles' awardee lists are public: fiscal 2023
at
\url{https://austin.widen.net/content/jfxwxsxktx/pdf/2023LiveMusicFundEventProgramAwardeeList.pdf}
and fiscal 2024 at
\url{https://austin.widen.net/content/bpzuacfkk8/pdf/FY24-AwardeeList_Austin-Live-Music-Fund.pdf}.
The applicant-level submissions the fund gathers, including geography,
sector and requested amount, have not been released in aggregate.}
& Consider publishing the applicant-level distributions the Live Music Fund
already gathers, in aggregate. This costs no new collection.
& The Texas Legislature attached a budget rider to the state film incentive
that requires the Governor's office to report film-incentive awards, recipient
by recipient, to the Legislative Budget Board (the Legislature's fiscal
research office) twice a year; no parallel rider covers the state music
rebate, and neither program is required to publish a recipient
list.\footnote{Rider 40 to the Governor's office in the 2026--27 General
Appropriations Act; Senate Bill 22, 89th Legislature, 2025. The parallel
music-rebate statute is Senate Bill 609, 87th Legislature, 2021.} \\
\rowrule
Between January 2020 and mid-2026, the Austin venue tracking panel recorded
more exits from its core live-music tier than it added, while the surviving
rooms in that tier traded no worse than the rest of the Austin mixed-beverage
market once venue-level trend and month-of-year seasonality are absorbed
(Section~\ref{sec:venues}).\footnote{The tracking panel is
\NVenuePanelIsCurated\ curated Austin live-music rooms matched to Texas
Comptroller mixed-beverage receipts (Open Data Portal dataset
\texttt{naix-2893} at
\url{https://data.texas.gov/dataset/Mixed-Beverage-Gross-Receipts/naix-2893}),
not a census. Section~\ref{sec:venues} tabulates the exits and the two-way
fixed-effects regression that supports the same-as-market finding, and the
technical appendix in Section~\ref{sec:appendix} lists the four ways the panel
can and cannot be read.}
& Consider directing support at room survival and replacement rather than at
operating margins. Property costs and lease security are among the pressures
venue operators report most often.
& Austin's own Iconic Venue Fund is aimed at exactly this. At its first award
in August 2023, requests ran to \NCityIvfRequestedToAwardRatio\ the size of
that award; the fund has awarded more since, and no updated recipient list
has been published (Section~\ref{sec:city}). \\
\rowrule
In the news timeline behind the venue panel, as many dated Austin
live-music-room closures cite a property cause (rent increase, sale of
property, redevelopment or lease expiration) as cite the pandemic, and
presenters who answered the 2022 Greater Austin Music Census ranked property
tax first among nine business pressures.\footnote{The closure counts come
from a 71-row news timeline compiled by this study from \emph{Austin
Chronicle}, \emph{Austin American-Statesman}, KUT/KUTX and industry
reporting; 36 rows carry a closure year and a coded cause (13 property, 13
pandemic, 2 regulatory, 3 uncoded, 5 other), and the property share appears
in the ledger as \texttt{census\_closures\_property\_share}. The presenter
ranking runs over the 2022 Greater Austin Music Census presenter-and-promoter
subsample (respondents self-identified as booking or presenting live music);
the census is a self-selected online sample, so the ranking describes
presenters who answered rather than a probability sample of Austin
presenters. Sources: news timeline at
\texttt{01\_evidence/08\_venues\_ecosystem/venue\_timeline.csv}; census
microdata dataset \texttt{an3p-3yqx} at
\url{https://data.austintexas.gov/}.}
& Consider whether venue support is best delivered through property-side
instruments rather than through grants.
& Other states run property or occupancy relief for cultural uses, and Texas
has built comparable mechanisms for other industries. \\
\rowrule
A musician applying to the Live Music Fund faced a different average award
depending on which fiscal year they applied, and the fund made no awards in
its fiscal 2025 cycle (Section~\ref{sec:city}).
& Consider fixing award sizes and cycle timing in advance and publishing the
schedule. Predictability is a design feature for people with irregular income.
& The fund's own year-to-year variation is the evidence; no external model is
needed. \\
\rowrule
The Texas Music Incubator Rebate returns a share of the mixed-beverage and
sales taxes a qualifying venue or promoter has already remitted, so the
subsidy scales with a venue's alcohol sales, and the statute requires a
written artist contract but sets no floor on what the contract pays the
performer.\footnote{Senate Bill 609, 87th Legislature Regular Session, 2021,
creating Government Code Chapter 485, Subchapter C, at
\url{https://capitol.texas.gov/tlodocs/87R/billtext/pdf/SB00609F.pdf}. Program
page: Texas Music Office, Texas Music Incubator Rebate Program, at
\url{https://gov.texas.gov/music/tmir}. Neither the statute nor the program
guidelines specify a minimum artist compensation.}
& Consider whether public support intended to reach musicians should be
conditioned on documented artist payments.
& Georgia authorised a music tax credit at \$15 million a year and paid out
nothing over the program's five-year life, so a headline cap alone says
nothing about what reached anyone.\footnote{Georgia Department of Audits and
Accounts and Georgia Southern University CBAER, \emph{Georgia Music Tax
Credit: Economic and Fiscal Impact Analysis}, December 2023, at
\url{https://www.audits.ga.gov/ReportSearch/download/30448}.} \\
\rowrule
Among respondents to the 2022 Greater Austin Music Census who answered the
three-year-retention items, housing tenure and health insurance separate who
expects to still be in Austin in three years; primary sector and workspace
status do not (Section~\ref{sec:census}).\footnote{The 2022 Greater Austin
Music Census carries no design weights, so every share here describes
respondents to a self-selected online instrument and not the Austin music
workforce; no margin of error can be computed for any of them. Development,
including the logit that isolates the tenure and insurance effects, is in
Section~\ref{sec:census}.}
& Consider treating musician retention as a housing question, with
instruments aimed at renters.
& Austin's cost problem concentrates in housing rather than in
overall living costs: the Austin metro's Regional Price Parity, the Bureau of
Economic Analysis series that indexes local price levels against a national
100, is near 100 on all items and above 100 on the housing item alone
(Section~\ref{sec:costs}). Any generalisation of these shares to the Austin
music workforce would rest on the assumption that the self-selected sample
resembles that workforce on tenure and insurance, which no source can
verify. \\
\rowrule
Federal shuttered-venue relief went to venues, cinemas and organisations
rather than to individuals, and the largest city instrument for buildings
(the Iconic Venue Fund) is capital support.\footnote{U.S. Small Business
Administration, Shuttered Venue Operators Grant awards file, awards as of
5~July 2022, at
\url{https://data.sba.gov/sites/default/files/distribution/SBA-ODA-2022-09-001/awards-as-of-7-5-22.xlsx}.
The program was authorised by section~324 of the Economic Aid to
Hard-Hit Small Businesses, Nonprofits, and Venues Act of 2020 and its
eligibility categories are venue operators, promoters, theatrical producers,
performing-arts organisations, motion picture theatre operators and talent
representatives, none of which is an individual worker.}
& Consider what share of support should follow the worker rather than the
building, and measure it.
& The Live Music Fund is the main city channel reaching individuals, and it
is the smallest of the three named here. Health coverage among Austin
musicians, at \NPumsAustinMusHicov\ (\NPumsAustinMusHicovMoe), is close within
that margin to \NPumsAustinAllworkHicov\ for all Austin
workers.\footnote{Both figures are from the American Community Survey
2020--2024 five-year Public Use Microdata Sample for Texas (dataset
released 5 March 2026 at
\url{https://www2.census.gov/programs-surveys/acs/data/pums/2024/5-Year/}),
Musicians and Singers (occupation code 2752) against all employed Austin
workers aged 18--64 with any-coverage flag HICOV=1, PWGTP-weighted, margin of
error at 90 percent from the 80 replicate weights; the musician cell rests on
\NPumsAustinMusNPos\ unweighted person records. See
Section~\ref{sec:appendix} for the population and weighting.} Austin's
independent nonprofit health programs for musicians reach individuals
directly, though this study cannot attribute the coverage rate to them. \\
\bottomrule
\end{longtable}
}

\vspace{5pt}
{\footnotesize\textbf{Cautions on reading this table.} Two policy
interventions this study tested left no detectable change in venue receipts:
the City of Austin's 2017 extended-hours pilot on Red River (a temporary
2 a.m.\ closing time for a defined set of Red River rooms; the design
detected no effect and rules out receipts effects larger than roughly 14
percent) and the fiscal 2024 round of Live Music Fund venue grants
(Section~\ref{sec:venues}). A third test found that receipts near the new
downtown arena that opened in April 2022 moved no differently from receipts
far from it, which is a test of proximity to the arena rather than of the
arena's overall effect. None of that means those interventions failed at
what they were trying to do; alcohol receipts at a venue are the wrong
instrument for detecting effects on what performers are paid, and no public
Austin dataset currently measures those payments directly. The largest
intervention during the study period was also federal, one-off, and
directed at venues and cinemas, so a local plan that assumes similar money
will come again is building on an unusual year.}
